\documentclass[reprint,amsmath,amssymb,aps,twoside]{revtex4-2} % for editor use only

% For initial submission use these lines
%\documentclass[amsmath,amssymb,aps,endfloats,linenumbers]{revtex4-2}



% DO NOT CHANGE THINGS BELOW HERE
\usepackage{graphicx}
\usepackage{amsmath,amssymb,amsfonts}
\usepackage{dcolumn}
\usepackage{bm}
\usepackage{siunitx}
%\usepackage{tikz,pgfplots}
\sisetup{separate-uncertainty=true, multi-part-units=single}
\usepackage[colorlinks,allcolors=blue]{hyperref}
\usepackage{cleveref}
\crefname{equation}{}{}
\crefname{figure}{Fig.}{Figs.}
\crefname{table}{Table}{Tables}
\usepackage{svg}
\svgpath{{./figures}}




\hypersetup{%
pdftitle={Non-conservation of mechanical energy in spring-launching toy poppers}, 
pdfauthor={Matthew Butch, Omar Ahmadzada, Jacob Pawelek, and Timothy Chung}, 
}


\usepackage{fancyhdr}
\pagestyle{fancy}
\fancyhf{}
\fancyhead[RE,RO]{J S\&E \textbf{2}, xx--xx (2026)} % FOR EDITOR ONLY
\fancyhead[LO]{Chung \emph{et al}} % if many authors
\fancyhead[LE]{Non-conservation of mechanical energy in poppers} 
\fancyfoot[C]{\thepage}
\fancypagestyle{mytitlepage}{
\fancyhf{}
\fancyhead[C]{Journal of Science \& Engineering \textbf{2}, xx--xx (2026)} % FOR EDITOR ONLY
\fancyfoot[C]{\thepage}
}


\begin{document}
%\setcounter{page}{67} % FOR EDITOR USE ONLY
%\linenumbers % FOR EDITOR USE ONLY

\title{Non-conservation of mechanical energy in spring-launching toy poppers}
\author{Matthew Butch}
\author{Omar Ahmadzada}
\author{Jacob Pawelek}
\author{Timothy Chung}
\email{Contact author: 427tchungr@frhsd.com}
\affiliation{Science \& Engineering Magnet Program, \href{https://manalapan.frhsd.com/}{Manalapan High School}, Englishtown, NJ 07726 USA}
\date{\today}

\begin{abstract}
The law of conservation of mechanical energy states that, in an isolated system, mechanical energy remains constant as energy transforms between its kinetic and potential form without loss. This experiment investigates if energy is conserved in a \qty{0.0055}{\kilo\gram} spring-launching toy popper. When compressed, the toy stores elastic potential energy that is then converted to kinetic energy as the toy is released. The kinetic energy is then converted to gravitational potential energy as the toy ascends. A two-tailed paired $t$-test comparing the initial kinetic energy and gravitational potential energy at maximum height across ten repeated trials of the same popper yielded a statistically significant difference from zero. This result appears to suggest that mechanical energy was not fully conserved, but a more likely explanation is that we did not account for the presence of non-conservative forces such as air resistance during launch.
\end{abstract}
\keywords{conservation of energy, mechanical energy, kinetic energy, potential energy, popper}

\maketitle\thispagestyle{mytitlepage}






%The current column width is \the\columnwidth

\section{Introduction}
The energy of motion, first introduced to describe the collision of two moving objects, is known as kinetic energy \cite{tipler,barrons,openstax}. This directionless, or scalar, quantity, is expressed in terms of mass $m$ and velocity $v$ and is given by
\begin{equation}
%KE=1/2mv^2 (Eq. 1)
KE = \frac{1}{2} m v^2.
\label{eq:1}
\end{equation}
The energy associated with an object’s position within a conservative force field is known as its potential energy, which is given by the negative of the work exerted by the a force, or
\begin{equation}
%PE=-integral(f*dr) (Eq. 2)
PE = -\int \vec{F}\cdot d\vec{r}.
\label{eq:2}
\end{equation}
The gravitational force exerted on objects near Earth’s surface has a downward direction with a magnitude of $mg$. Here, we’ll assume that upward vertical displacement $y$ is positive, so total displacement comes from initial point $y=0$ to final position $y=h$. We use this information to calculate the gravitational potential energy of an object, given by
\begin{equation}
GPE = - \int_0^h (-mg) dy,
\label{eq:3}
\end{equation}    
which evaluates to
\begin{equation}
GPE = mgh. 
\label{eq:4}
\end{equation}

A system’s total mechanical energy $E$ is given by the sum of its kinetic energy and potential energy. The law of energy conservation states that this sum will remain constant unless external forces do work on the system \cite{tipler,barrons,openstax}. Based on this law, we set up a vertical launch system where energy changes between different forms as some object rises against gravity. In this system, the initial kinetic energy of an object as it rises should be equivalent to its potential energy at the time the object reaches its maximum height and comes to a brief rest. Realistically, mechanical energy may not appear fully conserved in this system due to systemic deficits like air resistance and energy dissipation into heat and sound which may be unaccounted for. These factors arise since air molecules apply negative work on the system and the popper is not perfectly elastic \cite{collegephysics}. 

While these system energy losses cannot be accounted for with replication, we took ten trials of the same popper and evaluated the difference between energies in each trial. We establish the following:
%H0: the mean difference will be zero, consistent with energy conservation (mu d = 0)
%Ha: the mean difference will not be zero, inconsistent with energy conservation (mu d !=0)    







    

\section{Methods and materials}\label{sec:methods}

To begin our investigation on energy conservation, we first measured the mass of a single toy popper (Liberty Imports, USA) to be \qty{0.0055}{\kilo\gram} using a balance scale (Arboleaf, Model CK10l). Then, we stacked two meter sticks and pinned them against a wall so that we could properly measure up to two meters high. We also set up an iPhone 13 (Apple Inc; Cupertino, CA) parallel to the rulers to record our experimentation at 60FPS. We began collecting data by compressing the toy roughly \qty{0.025}{\meter}, as further compression would result in the toy snapping through, and then releasing it. After, we replayed the video to capture the greatest height of the toy and enter it into a table. This process was repeated for ten trials to account for random variation.

We also collected data for the force applied on the toy as its compression displacement increased. First, we placed a toy onto a balance and then compressed the toy in increments of \qty{0.005}{\meter}, up until the displacement is \qty{0.035}{\meter}. We chose this displacement amount since it would allow us to visualize the drop in force applied on the toy after the snap-through region of the toy is passed. This process was repeated twice because the force-displacement relationship in the popper is relatively consistent and not subject to random variation. For each trial, the mass displayed on the balance would be converted to Newtons and recorded in a force-displacement table.

We used \cref{eq:4} to calculate the gravitational potential energy at the peak height of the toy in each of our ten trials. We calculated the launch kinetic energy (LKE) across each ten trials by taking the data from our force-displacement table and creating a force-displacement function, for which we calculated the area under the curve up to the toy’s snap-through region. This area is equivalent to the toy’s elastic potential energy (EPE) in compression, which, in an energy-conserving system, is equal to the toy’s LKE. The math can be written as the following:
\begin{equation}
%KE at launch ≈ EPE ≈ integral(force*displacement) from displacement = 0 to = (displacement) (Eq. 5)
KE_{launch} \approx EPE \approx \int_0^{\Delta x} F dx.
\end{equation}

We evaluate the integral once for each ten trials using the displacement of compression in each
using a trapezoidal sum. Finally, using the Microsoft Excel addon Analysis ToolPak, we ran a two-tailed paired $t$-test to determine whether or not the mean difference in energies is zero to some degree of uncertainty. Data and analysis code are provided at \url{https://github.com/devangel77b/427tchung-lab3}. 




    
\section{Results}
\begin{table}
\caption{Two trials of measuring spring force (\unit{\newton}) and displacement (\unit{\meter})}
\label{tab:1}
\begin{center}
\begin{tabular}{cc}
displacement, \unit{\meter} & spring force, \unit{\newton} \\
\colrule
0 & 0 \\
0.005 & 1.442 \\
0.005 & 1.501 \\
0.010 & 2.816 \\
0.010 & 2.953 \\
0.015 & 4.239 \\
0.015 & 4.374 \\
0.020 & 5.778 \\
0.020 & 5.603 \\
0.025 & 7.770 \\
0.025 & 7.565 \\
0.030 & 6.710 \\
0.030 & 6.889 \\
0.035 & 5.374 \\
0.035 & 5.943 
\end{tabular}
\end{center}
\end{table}

\begin{figure}
\begin{center}
% PLACEHOLDER
\includesvg[width=\columnwidth]{fig1.svg}
\end{center}
\caption{$F$ vs. $x$ with two force data points per displacement indicating a replicate trial. Popper snap-through instability is visualized with a decrease in spring force around \qty{0.025}{\meter}.}
\label{fig:2}
\end{figure}

As described in \cref{sec:methods}, the area under the $F$ vs. $x$ curve was calculated for each compression displacement of the ten trials using a trapezoidal sum, which yielded the LKE values that were inputted into \cref{tab:2}.

\begin{table}
\caption{KE, PE, and the differences}
\label{tab:2}
\begin{center}
% PLACEHOLDER
\includegraphics[width=\columnwidth]{tables/table2.png}
\end{center}
\end{table}

\begin{table}
\caption{Mean and standard deviation of KE, PE, and KE-PE}
\label{tab:3}
\begin{center}
% PLACEHOLDER
\includegraphics[width=\columnwidth]{tables/statistics.png}
\end{center}
\end{table}

A two-tailed paired $t$-test using values from \cref{tab:3} yielded a $p=\num{3e-6}$.











\section{Discussion}

\subsection{Interpretation}
The goal of this experiment was to test the law of energy conservation using a toy popper. Overall, we were unable to accurately present this law. We wrote this lab with the intention of showing that the mean difference in LKE and GPE would not significantly differ from zero, since that would be an effective representation of the law. Our statistical test, however, yielded a $p$-value of $p=\num{3e-6}$, a value far below any common significance level, meaning that there is a significant difference between the LKE and GPE. Since the difference between LKE and GPE was positive across each trial, we have a potential indicator that there were non-conservative forces acting on the toy which resulted in energy dissipation as it ascended. Something like air resistance may have applied negative work on the system, dissipating energy and reducing the amount that can be converted into GPE \cite{tipler,barrons,openstax}.

\subsection{Sources of experimental error}
The human eye being unable to look at two places simultaneously may have been a potential confounding variable. When recording the height of the toy, we had to look at the toy scanning for when it would approximately have zero speed in the air. We would then quickly look at the meter stick taking that height to account. This height may have been overestimated or underestimated depending on when the person watching the toy looked at the connected meter sticks on the right. We took 10 trials and averaged this out on purpose to try and avoid this, however it most likely still affected our results. Another potential source of error is ignoring air drag, which likely went against the object’s motion; failure to account for this may provide an underestimate of the total energy in the system.






    
    
\section{Acknowledgements}
We thank the Science and Engineering faculty of Manalapan High School for providing us the means to conduct our experimental research. JP was responsible for data collection, OA wrote the abstract and interpreted the results, MB wrote the lab setup and created diagrams to represent it, TC wrote the introduction and the results of the experiment.






\bibliography{lab.bib}
%References
%[1]“Conservation of Energy - University Physics Volume 1 - OpenStax,” Openstax.org, 2016.
%https://openstax.org/books/university-physics-volume-1/pages/8-3-conservation-of-energy
%[2]“7.2 Kinetic Energy - University Physics Volume 1 | OpenStax,” openstax.org.
%https://openstax.org/books/university-physics-volume-1/pages/7-2-kinetic-energy
%[3]“2.7 Falling Objects - College Physics 2e | OpenStax,” openstax.org.
%https://openstax.org/books/college-physics-2e/pages/2-7-falling-objects

\end{document}
